\subsubsection{Material set models and categorical CH}
\label{sec:p10-set-models-structural-ch}

A material model carries a membership relation. A topos instead
carries objects and morphisms. A comparison between their continuum
statements must specify how one structure is interpreted in the other.

\begin{hbdefinition}{def:p10-set-model-category}{The category coded by a set model}
Let \(M\) be a set-sized model of \(\mathsf{ZFC}\), not necessarily
well-founded or transitive.
The objects of \(\mathcal C_M\) are the elements of its domain.
An arrow from \(x\) to \(y\) is an \(M\)-coded function that \(M\)
regards as having domain \(x\) and values in \(y\), together with
those source and target objects. Identities and composition are
the operations computed in \(M\).
All membership and function assertions in this definition use
the relation of \(M\), not ambient membership.
\end{hbdefinition}

\begin{hbtheorem}{thm:p10-set-model-topos}{proved=annex:p10-set-model-topos}{The structural set model}
The category \(\mathcal C_M\) is a nondegenerate, well-pointed
Boolean topos with a natural numbers object, and every epimorphism
splits. Its finite limits, finite colimits, exponentials, and
subobject classifier are the constructions coded in \(M\).
Its natural numbers object is \(M\)'s \(\omega\).
The construction does not identify its arrows with all external
functions between externally interpreted sets.
\end{hbtheorem}

For objects \(X,Y\), use the source's strict comparison:
there is a monomorphism \(X\to Y\), and the internal object of
epimorphisms in that direction is empty.
The categorical continuum assertion says that no object lies
strictly between \(\omega\) and \(2^\omega\).
In the following comparison all of these objects and predicates
belong to \(\mathcal C_M\).

\begin{hbtheorem}{thm:p10-set-model-ch}{proved=annex:p10-set-model-ch}{Exact comparison of continuum statements}
The categorical continuum assertion holds in \(\mathcal C_M\)
if and only if \(M\models\mathsf{CH}\).
Both the exponential \(2^\omega\) and its epimorphism predicates
are computed in \(M\), not in the external metatheory.
\end{hbtheorem}

\begin{hbcorollary}{thm:p10-structural-ch-consistency}{proved=annex:p10-structural-ch-consistency}{The two consistency comparisons}
If \(\mathsf{ZF}\) is consistent, the structural set axioms
listed in \cref{thm:p10-set-model-topos} have models satisfying
the categorical continuum assertion and models satisfying its
negation.
The positive material consistency argument is supplied by
\cref{thm:p06-gch-inner-l,thm:ch-relative-consistency};
the negative material argument is the separate Boolean-valued
construction of \cref{thm:p06-bfc-consistency}.
\end{hbcorollary}

This interpretation runs from material models to categorical
structure. It does not reconstruct a membership universe from an
arbitrary topos or infer material Replacement from elementary
topos axioms. In particular, the material non-CH result cited here
uses its own Boolean-valued membership-model proof; it is not
inferred merely from the internal strict inequalities in the
sheaf construction.
